\centerline{\Huge \bf  Тренировочная Олимпиада 1}
\centerline{\Huge \bf 8 класс}
\centerline{\Huge \bf Уровень: городская Олимпиада.}


\begin{flushright}
    \scriptsize{Глава 1. Начало конца.}
\end{flushright}

\problem{1} 
Татары считают, что среднее количество поедаемого чак-чака их населением в 41 раз превосходит среднее количество поедаемого чак-чака всеми жителями России (включая всех татар). Докажите, что татары лукавят. Население татарстана составляет 2,5\% от населения России.
\begin{flushright}
    \textit{\small{(Гасанов Р.В.)}}
\end{flushright}

\problem{2} За победу в кубке ЧГК ЭлМШ команда Германа Эрдема Олега выиграли \\987654321123456789281 конфет. Но оказалось, что они не способны поровну разделить конфеты между собой. Анджур Аркадьевич как ответственный организатор, чтоб дети не поссорились решил незаметно для всех съесть несколько конфет. Какое максимальное количество конфет может достаться каждому победителю?
\begin{flushright}
    \textit{\small{(Очиров О.В.)}}
\end{flushright}

\problem{3} $P(x), Q(x), R(x) \ -$ различные приведенные квадратные трёхчлены, с положительным дискриминантом. Какое наибольшее количество корней может иметь уравнение $(P(x)Q(x))^2+(Q(x)R(x))^2+(R(x)P(x))^2 = 0$
\begin{flushright}
    \textit{\small{(Гасанов Р.В.)}}
\end{flushright}

\problem{4}  Найти коэффициент при $x^{48}$ у многочлена $(1 + x^8 + x^{12} + x^{24})^6$

\begin{flushright}
    \textit{\small{(Гасанов Р.В.)}}
\end{flushright}


\problem{5} В остроугольном треугольнике $MNK$ продлили сторону MK за точку K на отрезок $KE$ такой, что $\dfrac{KE}{MK} = \dfrac{3}{25}$. Биссектриса $MQ$ $-$ треугольника $MNE$ пересекает $NK$ в точке $P$ так, что $NP = PK$. Найдите отношение площади треугольника $MNP$ к площади четырёхугольника $PQEK$.

\begin{flushright}
    \textit{\small{(Доржеев А.А.)}}
\end{flushright}